% TOP_PROBLEM: {"problemId": "problem.local-smoothing-conjecture-for-wave-equations", "canonicalTitle": "Local Smoothing Conjecture for Wave Equations", "releaseRank": 201, "primaryDomain": "analysis_pde", "primaryDomainLabel": "Analysis & PDE", "displayStatus": "open", "releaseStatus": "open"}
% !TeX program = lualatex
% Definitions and sources retained from research_batch_201_202.tex; research is in GPT_6_Astra_Ultra.
\documentclass[11pt]{article}
\usepackage[margin=25mm]{geometry}
\usepackage{fontspec}
\setmainfont{Latin Modern Roman}
\usepackage{amsmath,amssymb,mathtools}
\usepackage{unicode-math}
\setmathfont{Latin Modern Math}
\usepackage{enumitem,needspace}
\usepackage{xurl,hyperref}
\hypersetup{colorlinks=true,urlcolor=blue}
\setcounter{secnumdepth}{1}
\setlength{\parindent}{0pt}
\setlength{\parskip}{5pt}
\setlength{\emergencystretch}{3em}
\setlist[itemize]{leftmargin=3em,itemsep=5pt,topsep=4pt,parsep=0pt}
\allowdisplaybreaks
\newcommand{\N}{\mathbb{N}}
\newcommand{\Z}{\mathbb{Z}}
\newcommand{\Q}{\mathbb{Q}}
\newcommand{\R}{\mathbb{R}}
\newcommand{\C}{\mathbb{C}}
\newcommand{\F}{\mathbb{F}}
\newcommand{\A}{\mathbb{A}}
\newcommand{\PP}{\mathbb P}
\newcommand{\E}{\mathbb{E}}
\newcommand{\eps}{\varepsilon}
\newcommand{\dd}{\,\mathrm d}
\newcommand{\sref}[2]{\hyperlink{source:#1:#2}{[S#2]}}
\newcommand{\eref}[2]{\hyperlink{extra:#1:#2}{[E#2]}}
% Turn the next switch off for a shorter reading copy. The quotations preserve
% every exactTarget field, including qualifications omitted from a short summary.
\newif\ifshowcatalogue
\showcataloguetrue
\newcommand{\cataloguescope}[1]{\ifshowcatalogue
	\paragraph{Exact scope in the supplied catalogue.}
	\begin{quote}\small #1\end{quote}\fi}
\begin{document}
\setcounter{section}{200}
	% ================================================================
	% problemId: problem.local-smoothing-conjecture-for-wave-equations
	% source releaseRank: 201; source releaseStatus: open
	% exactTarget SHA-256: 69808b9a81b2b9968148156e4e2a52f3bd538cfa32b8ff88f04e4a834f21f1b0
	\clearpage
	\section{\texorpdfstring{Local Smoothing Conjecture for Wave Equations}{Local Smoothing Conjecture for Wave Equations}}\label{201}
	\subsection{Definitions and mathematical statement}
	For a Schwartz function on $\R^n$, the half-wave propagator is the Fourier multiplier
	$\widehat{e^{it\sqrt{-\Delta}}f}(\xi)=e^{it|\xi|}\widehat f(\xi)$, using the Fourier convention in which $-\Delta$ has symbol $|\xi|^2$. Define the Bessel-potential norm
	$\|f\|_{W^{s,p}}=\|(1-\Delta)^{s/2}f\|_{L^p}$ and put $s_p=(n-1)(1/2-1/p)$. The assertion is
	\[
	\left(\int_1^2\int_{\R^n}|e^{it\sqrt{-\Delta}}f(x)|^p\dd x\dd t\right)^{1/p}
	\le C_{n,p,\eps}\|f\|_{W^{s_p-1/p+\eps,p}}
	\]
	for $n\ge3$, $p>2n/(n-1)$, and every $\eps>0$. Constants are independent of $f$. This is a spacetime estimate with an arbitrarily small derivative loss; the excluded endpoint and a zero-loss bound are not asserted.
	\par\smallskip\textit{Formulation and concept references:} \sref{201}{1}.
	\cataloguescope{For every n >= 3, p > 2n/(n-1), and epsilon > 0, prove ||exp(it sqrt(-Delta))f||\_\{L\textasciicircum{}p(R\textasciicircum{}n x [1,2])\} <= C\_\{p,epsilon\} ||f||\_\{W\textasciicircum{}\{s\_p-1/p+epsilon,p\}(R\textasciicircum{}n)\}, where s\_p=(n-1)(1/2-1/p).}
	\subsection{Short English statement}
	Does averaging the wave evolution over time save almost one-over-p derivatives throughout the conjectured range of exponents?
	\subsection{Sources}
	\begin{itemize}\raggedright
		\item[\textnormal{[S1]}]\hypertarget{source:201:1}{} Formal statement, Introduction, local smoothing conjecture.
		\url{https://www.sciencedirect.com/science/article/pii/S0022123623003786}.
		{\small\textit{Catalogue formulation source.}}
		% BEGIN RESEARCH REFERENCES 201
		\item[\textnormal{[E1]}]\hypertarget{extra:201:1}{}
		J. Bourgain and C. Demeter, \emph{The proof of the $\ell^2$ Decoupling
			Conjecture}, Annals of Mathematics (2) 182 (2015), no.~1, 351--389,
		Theorem 1.2 (cone decoupling) and the discussion immediately following it.
		\url{https://annals.math.princeton.edu/2015/182-1/p09};
		author manuscript \url{https://arxiv.org/pdf/1403.5335v3}.
		{\small\textit{Published result; consulted 24 September 2026.}}
		\item[\textnormal{[E2]}]\hypertarget{extra:201:2}{}
		S. Gan, D. He, X. Li, and S. Wu, \emph{On local smoothing estimates for
			wave equations}, arXiv:2502.05973v2, 4 January 2026, Theorem 1.3,
		equation (1.4), and Sections 1.1, 1.3, and 5.1.
		\url{https://arxiv.org/html/2502.05973v2}.
		{\small\textit{Recent manuscript theorem, not independently audited here;
				version 2 has four authors and an improved main result. Consulted
				24 September 2026.}}
		\item[\textnormal{[E3]}]\hypertarget{extra:201:3}{}
		D. Beltran, J. Roos, A. Rutar, and A. Seeger, \emph{A fractal local
			smoothing problem for the wave equation}, Bulletin of the London
		Mathematical Society 57 (2025), no.~12, 3667--3690; author manuscript
		arXiv:2501.12805v1, 22 January 2025, Theorem 1.1, equation (1.6),
		and Section 5 (radial inputs).
		\url{https://arxiv.org/html/2501.12805v1}.
		{\small\textit{The radial theorem, not its January 2025 general-frontier
				summary, is used here. Consulted 24 September 2026.}}
		% END RESEARCH REFERENCES 201
	\end{itemize}
	
\end{document}
